\section{Methods}
\label{sec:methods}

This section describes the full experimental design: the datasets, the
construction that holds training-set size fixed across imbalance ratios, the
two directed label-error mechanisms, the seven correction methods audited, the
metrics used to score them, and the statistical tests used to compare them.

\subsection{Overview}
\label{sec:methods-overview}

Figure~\ref{fig:pipeline} summarizes the pipeline applied to every dataset. A
dataset is split once into a training pool and a held-out test set, which stays
fixed and clean for the rest of that seed's run. Only the training side is
resampled and corrupted, and each correction method is fit on that side and
scored against the same held-out test set. A controlled sample is drawn
from the training pool at a target minority ratio $r$ and a fixed total size
$N_{\text{common}}$ (Section~\ref{sec:methods-ncommon}). A directed label error
is then injected into that sample at a contamination rate $c$, in one of two
directions (Section~\ref{sec:methods-corruption}), or left clean at $c{=}0$.
Every correction method (Section~\ref{sec:methods-corrections}) is fit on the
resulting training set and scored against the same held-out test set, and the
scores are converted into the metrics described in
Sections~\ref{sec:methods-metrics} and~\ref{sec:methods-contribution}. This is
repeated for every combination of dataset, seed, imbalance ratio, contamination
rate, and direction, and for every correction method, producing one row per
method per condition.

\begin{figure*}[t]
\centering
\resizebox{\linewidth}{!}{%
\begin{tikzpicture}[
  node distance=0.35cm,
  box/.style={draw, rounded corners, align=center, minimum height=1.15cm,
              text width=1.85cm, font=\footnotesize},
  arr/.style={-{Stealth[length=2mm]}, thick}
]
\node[box] (data) {Dataset\\(1 of 26 UCI\\sources)};
\node[box, right=of data] (split) {Stratified\\80/20 split\\(test set\\held out)};
\node[box, right=of split] (sample) {Controlled\\sample\\($N_{\text{common}}$,\\ratio $r$)};
\node[box, right=of sample] (corrupt) {Inject label\\error ($c$,\\Direction A/B,\\or clean)};
\node[box, right=of corrupt] (fit) {Fit\\correction\\method};
\node[box, right=of fit] (predict) {Predict on\\held-out\\test set};
\node[box, right=of predict] (score) {Metrics \&\\CG / CHR /\\$c^*$ / PSG};

\draw[arr] (data) -- (split);
\draw[arr] (split) -- (sample);
\draw[arr] (sample) -- (corrupt);
\draw[arr] (corrupt) -- (fit);
\draw[arr] (fit) -- (predict);
\draw[arr] (predict) -- (score);
\end{tikzpicture}%
}
\caption{Stress-cell pipeline, applied to each dataset, seed, ratio,
contamination rate, and direction.}
\label{fig:pipeline}
\end{figure*}

\subsection{Datasets}
\label{sec:methods-datasets}

Twenty-six binary-classification datasets were drawn from the UCI Machine
Learning Repository~\cite{kelly_uci}: Banknote
Authentication~\cite{lohweg2012banknote},
Blood Transfusion~\cite{yeh2008bloodtransfusion}, South German
Credit~\cite{groemping2019germancredit}, Spambase~\cite{hopkins1999spambase},
Breast Cancer Wisconsin~\cite{wolberg1993breastcancer},
Ionosphere~\cite{sigillito1989ionosphere}, Haberman's
Survival~\cite{haberman1999survival}, Statlog
Heart~\cite{statlogheart}, HTRU2~\cite{lyon2015htru2}, QSAR
Biodegradation~\cite{mansouri2013qsar}, Adult~\cite{becker1996adult},
Bank Marketing~\cite{moro2014bankmarketing}, MAGIC Gamma
Telescope~\cite{bock2004magic}, Online Shoppers Purchasing
Intention~\cite{sakar2018onlineshoppers}, Congressional Voting
Records~\cite{uci1987congressionalvoting}, Chronic Kidney
Disease~\cite{rubini2015ckd}, Default of Credit Card
Clients~\cite{yeh2009creditcard}, Taiwanese Bankruptcy
Prediction~\cite{uci2020taiwanbankruptcy}, Polish Companies
Bankruptcy~\cite{tomczak2016polishbankruptcy}, Mushroom~\cite{uci1981mushroom},
Cervical Cancer (Risk Factors)~\cite{fernandes2017cervical},
SECOM~\cite{mccann2008secom}, Arrhythmia~\cite{guvenir1997arrhythmia},
Internet Advertisements~\cite{kushmerick1999internetads},
Seismic-Bumps~\cite{sikora2010seismic}, and APS Failure at Scania
Trucks~\cite{scania2016apsfailure}, chosen so that no single result could be
attributed to one dataset's idiosyncrasies, spanning finance, security,
healthcare, signal processing, astronomy, chemistry, industrial
manufacturing, retail behavior, and political science. Each dataset
designates one binary class as its positive, or target, class
(Table~\ref{tab:datasets}); in the controlled training sets built in
Section~\ref{sec:methods-ncommon}, this class's prevalence is fixed to
$r$, so it is the minority class at every ratio below $r{=}0.50$ and
exactly balanced at $r{=}0.50$, regardless of how common or rare it is
in the original, unmodified dataset. Together they range
from 270 to 60,000 raw rows and from 3 to 1{,}558 raw features, and their
natural positive-class prevalence ranges from $1.7\%$ (APS Failure) to
$62.5\%$ (Chronic Kidney Disease, the one dataset in this study where the
positive class is the natural majority). Table~\ref{tab:datasets}
summarizes all 26 datasets' size, the controlled training size
$N_{\text{common}}$ every stress cell actually draws from, and the meaning
of the positive class in each. The $N_{\text{common}}$ column is the
controlled training size (Section~\ref{sec:methods-ncommon}) that every stress
cell draws from, not the raw row count. It is computed at seed 13, and
seed-to-seed variation is at most a few rows from stratified-split rounding.
Feature counts marked $\to$100 are capped (Section~\ref{sec:methods-datasets}).

\begin{table*}[t]
\centering
\caption{Dataset characteristics for all 26 datasets.}
\label{tab:datasets}
\begin{tabular}{@{}lrrrl@{}}
\toprule
Dataset & Rows & Feat. & $N_{\text{common}}$ & Positive class \\
\midrule
Banknote Authentication      & 1{,}372 & 4  & 641   & Forged banknote \\
Blood Transfusion             & 748    & 4  & 284   & Blood donated in target period \\
South German Credit           & 1{,}000 & 20 & 480   & Bad credit risk \\
Spambase                      & 4{,}601 & 57 & 2{,}347 & Spam email \\
Breast Cancer Wisconsin       & 569    & 30 & 300   & Malignant diagnosis \\
Ionosphere                    & 351    & 34 & 189   & No ionospheric structure \\
Haberman's Survival           & 306    & 3  & 128   & Death within 5 years \\
Statlog Heart                 & 270    & 13 & 126   & Heart disease present \\
HTRU2                         & 17{,}898 & 8 & 2{,}622 & Pulsar candidate \\
QSAR Biodegradation           & 1{,}055 & 41 & 568   & Ready biodegradable \\
Adult                          & 48{,}842 & 14 & 18{,}698 & Income exceeds \$50K/year \\
Bank Marketing                 & 45{,}211 & 16 & 8{,}462  & Client subscribed to term deposit \\
MAGIC Gamma Telescope          & 19{,}020 & 10 & 10{,}385 & Hadron (background) event \\
Online Shoppers                & 12{,}330 & 17 & 3{,}052  & Session ended in completed purchase \\
Congressional Voting           & 435     & 16 & 225    & Republican representative \\
Chronic Kidney Disease         & 400     & 24 & 126    & Chronic kidney disease present \\
Default of Credit Card Clients & 30{,}000 & 23 & 10{,}618 & Default on payment next month \\
Taiwanese Bankruptcy           & 6{,}819  & 95 & 352    & Company went bankrupt \\
Polish Companies Bankruptcy    & 43{,}405 & 65 & 3{,}346  & Bankrupt within forecast period \\
Mushroom                       & 8{,}124  & 22 & 3{,}543  & Poisonous mushroom \\
Cervical Cancer (Risk Factors) & 858     & 32 & 88     & Biopsy-confirmed abnormality \\
SECOM                          & 1{,}567  & 590$\to$100 & 166 & Semiconductor test fail \\
Arrhythmia                     & 452     & 279$\to$100 & 206 & Cardiac arrhythmia present \\
Internet Advertisements        & 3{,}279  & 1{,}558$\to$100 & 734 & Image is an advertisement \\
Seismic-Bumps                  & 2{,}584  & 18 & 272    & Hazardous seismic state next shift \\
APS Failure at Scania Trucks   & 60{,}000 & 170$\to$100 & 1{,}600 & APS component failure \\
\bottomrule
\end{tabular}
\end{table*}

No transformation altered the semantics of any feature. Two categories of
transformation changed feature values or feature counts, neither altering
what a feature means. First, column renames, purely cosmetic, applied to
Spambase and QSAR Biodegradation: XGBoost's tree serialization format
rejects the characters \texttt{[}, \texttt{]}, and \texttt{<} in feature
names, so Spambase's \texttt{char\_freq\_[} (the frequency of the
\texttt{[} character in an email) was renamed to
\texttt{char\_freq\_bracket\_open}, and eight of QSAR Biodegradation's
molecular-descriptor columns (e.g., \texttt{F01[N-N]}) had their brackets
replaced with a hyphen (e.g., \texttt{F01\_N-N}), with the underlying values
unchanged in both cases. Second, a feature cap applied to the four datasets
with more than 100 raw features, SECOM (590), Arrhythmia (279), Internet
Advertisements (1{,}558), and APS Failure (170): each is capped to a
deterministic random subsample of 100 columns (fixed seed, sorted for a
stable, reviewable column order), so that no method's training cost scales
with raw feature count rather than with the imbalance and noise conditions
actually under test. The columns kept are fixed once per dataset and
identical across every seed, ratio, and contamination level tested on it.

The pipeline separately includes leakage-safe median imputation for numeric
columns and most-frequent-value imputation followed by one-hot encoding for
categorical columns, fit only on each stress condition's training rows.
Ten of the 26 datasets, Banknote Authentication, Blood Transfusion, South
German Credit, Spambase, Breast Cancer Wisconsin, Ionosphere, Haberman's
Survival, Statlog Heart, HTRU2, and QSAR Biodegradation, contain no missing
values and no non-numeric feature column (South German Credit's twenty
attributes and Statlog Heart's categorical clinical measures are
integer-coded and are treated as numeric throughout), so neither step
alters any value for those ten. Of the remaining sixteen, eleven contain at
least one column with missing values (Adult, Bank Marketing, Congressional
Voting, Chronic Kidney Disease, Polish Companies Bankruptcy, Mushroom,
Cervical Cancer, SECOM, Arrhythmia, Internet Advertisements, and APS
Failure) and seven contain genuinely categorical, non-numeric columns
(Adult, Bank Marketing, Online Shoppers, Congressional Voting, Chronic
Kidney Disease, Mushroom, and Seismic-Bumps), so both imputation steps and
the one-hot encoding step are genuinely exercised in this study, rather
than being present in the pipeline but never triggered.

Before any sampling or corruption, each dataset is split once with a
stratified 80/20 train/test split, so both partitions preserve the dataset's
original class ratio. The test partition is fixed for that seed and never
altered by any later step. Every correction method, at every imbalance ratio
and contamination level, is scored against exactly the same test rows.

That fixed test partition has a consequence worth naming directly. Because
it is never resampled, it keeps each dataset's natural positive-class
prevalence throughout, while Section~\ref{sec:methods-ncommon} deliberately reshapes
the training partition to $r$. Natural prevalence across the 26 datasets
ranges from $1.7\%$ (APS Failure) to $62.5\%$ (Chronic Kidney Disease).
Seven datasets sit at or below the two most severe tested ratios: APS
Failure ($1.7\%$) and Taiwanese Bankruptcy ($3.2\%$) fall below even
$r{=}0.05$; Polish Companies Bankruptcy ($4.8\%$) sits almost exactly on
it; and Cervical Cancer ($6.4\%$), Seismic-Bumps ($6.6\%$), SECOM
($6.6\%$), and HTRU2 ($9.2\%$) all sit close to $r{=}0.10$. So at
$r{=}0.05$ and $r{=}0.10$ especially, training and test are drawn from
different class priors for a substantial share of the datasets in this
study. The severe-imbalance conditions this study emphasizes are therefore
not imbalance in isolation. They combine
imbalance with a train/test prior shift, exactly the kind of mismatch
DistPFN and DistPFN-T are built to correct
(Section~\ref{sec:methods-corrections}), and the kind
\texttt{tabpfn\_threshold}'s own threshold adjustment addresses as well.
Section~\ref{sec:discussion-simplicity} returns to what this means for
interpreting why the simplest corrections in this study lead, and
Section~\ref{sec:threats-internal} states the resulting limitation on
treating imbalance and prior shift as separable.

\subsection{Controlled Imbalance: A Common Training Size}
\label{sec:methods-ncommon}

Each dataset is tested at four minority ratios, $r \in \{0.50, 0.20, 0.10,
0.05\}$, where $r$ is the fraction of training rows that belong to the
minority class. Comparing a model trained at a high ratio (e.g., $r{=}0.20$)
against one trained at a low ratio (e.g., $r{=}0.05$) is only informative
about imbalance itself if both training sets are the same total size.
Otherwise, a difference in performance could be explained by sample size
rather than by imbalance. To remove that confound, one shared training size
$N_{\text{common}}$ is computed for each dataset and seed and used at every
tested ratio. Only the split between minority and majority rows within that
fixed total changes as $r$ changes.

Let $n_{\min}$ and $n_{\max}$ denote the number of positive-class and
other-class rows available in the training pool for a given dataset and
seed (named for their role in the controlled sample below, not for which
is numerically larger in the raw pool; Chronic Kidney Disease's positive
class is the raw pool's natural majority, Section~\ref{sec:methods-datasets}).
At ratio $r$, a
total training size of $N$ requires $rN$ minority rows and $(1-r)N$ majority
rows, so the largest $N$ achievable at that ratio is bounded by whichever
class runs out first:
\begin{equation}
\text{feasible\_n}(r) = \min\!\left(\frac{n_{\min}}{r},\ \frac{n_{\max}}{1-r}\right).
\label{eq:feasible-n}
\end{equation}
$N_{\text{common}}$ is the largest size achievable at every tested ratio
simultaneously, i.e., the tightest of these four constraints, floored to a
whole number of rows:
\begin{equation}
N_{\text{common}} = \left\lfloor \min_{r \in \{0.50,0.20,0.10,0.05\}} \text{feasible\_n}(r) \right\rfloor.
\label{eq:ncommon}
\end{equation}
For a given $r$ and $N_{\text{common}}$, a controlled training set is drawn by
sampling $\text{round}(r \cdot N_{\text{common}})$ minority rows and the
remaining $N_{\text{common}}$ rows from the majority class, both without
replacement, then shuffling the combined set.

\subsection{Directed Label-Error Mechanisms}
\label{sec:methods-corruption}

Two label-error mechanisms are applied to the controlled training set built
in Section~\ref{sec:methods-ncommon}, each at contamination rate $c \in
\{0.05, 0.10, 0.20, 0.30\}$, plus a shared clean condition at $c{=}0$. Both
flip a controlled number of true labels within that set, but they differ in
which class is flipped and how $c$ is defined relative to the resulting
pool. Here $n_{\min}$ denotes the number of true minority rows in that
controlled training set, equal to $\text{round}(r \cdot N_{\text{common}})$
by construction, not the pool-level $n_{\min}$ from
Section~\ref{sec:methods-ncommon}.

\textbf{Direction A (contamination): majority rows relabeled as minority.}
$c$ is defined as the fraction of the \emph{observed} minority-labeled pool
that is actually mislabeled majority rows. If $k$ majority rows are
flipped,
\begin{equation}
c = \frac{k}{n_{\min} + k}
\quad\Longrightarrow\quad
k = \frac{c}{1-c}\, n_{\min}.
\label{eq:directionA}
\end{equation}
$k$ rows are drawn uniformly at random from the majority pool and relabeled.

\textbf{Direction B (loss): minority rows relabeled as majority.}
Here $c$ is defined directly against the true minority count:
\begin{equation}
k = c \cdot n_{\min}.
\label{eq:directionB}
\end{equation}
$k$ rows are drawn uniformly at random from the minority pool and relabeled.

Both mechanisms return the corrupted labels together with the realized
contamination rate (which can differ slightly from the requested $c$ due to
integer rounding of $k$) and the indices flipped, so every corruption is
fully auditable. This asymmetric, class-conditional construction follows the
label-noise formalization of Natarajan et al.~\cite{natarajan2013noisylabels}
and Scott et al.~\cite{scott2013asymmetricnoise}, rather than symmetric random
flipping.

That difference in definition is not cosmetic. Direction A's $c$ describes
contamination as it would be observed after the fact, the fraction of
everything now labeled minority that is actually a mislabeled majority row,
close to asking how many flagged cases turn out to be false alarms.
Direction B's $c$ instead measures signal loss directly against the true
minority count, the fraction of genuine minority rows destroyed. The two are
natural for the process each direction represents, but they are not the
same quantity. The same nominal $c$ therefore corresponds to a different
flip count relative to $n_{\min}$ in each direction. At $c{=}0.30$,
Direction A flips $k/n_{\min} = c/(1-c) = 42.9\%$ of $n_{\min}$'s worth of
majority rows, while Direction B flips exactly $30\%$, a gap that widens
with $c$ (from $5.3\%$ against $5.0\%$ at $c{=}0.05$) but stays independent
of $r$, since both ratios are taken relative to $n_{\min}$ within their own
direction's formula rather than against the majority count. Comparing
directions at matched nominal $c$, as Section~\ref{sec:results-rq2} does,
is therefore a comparison of each direction's own natural severity, not of
an identical absolute flip count, a distinction
Section~\ref{sec:threats-internal} returns to.

\subsection{Correction Methods}
\label{sec:methods-corrections}

Seven correction methods are audited, split into two families so that every
method is compared against the uncorrected baseline for the same underlying
model. All XGBoost variants share one frozen configuration
(Table~\ref{tab:config}), never tuned per condition, so any performance
difference between them is attributable to the correction technique itself.

\subsubsection{XGBoost family}

\texttt{xgb\_none} trains XGBoost~\cite{chen2016xgboost} directly on the
corrupted, imbalanced training set with no correction, and is the baseline
every other XGBoost variant is measured against. \texttt{xgb\_weighted} sets
\texttt{scale\_pos\_weight} to the ratio of observed majority to observed
minority rows, following the cost-sensitive framing of
Elkan~\cite{elkan2001costsensitive}. \texttt{xgb\_ros} applies random
oversampling, duplicating minority-labeled rows with replacement until the
classes are balanced. \texttt{xgb\_smote} instead synthesizes new
minority-labeled rows with SMOTE~\cite{chawla2002smote}, interpolating
between a minority row and one of its $k$ nearest minority neighbors, a
parameter unrelated to the flip count $k$ in
Section~\ref{sec:methods-corruption}. It is set to
$k = \min(5,\, n_{\min}^{\text{obs}}-1)$, where $n_{\min}^{\text{obs}}$ is the
observed minority count in the corrupted training set actually being
resampled (distinct from both $n_{\min}$ in
Section~\ref{sec:methods-ncommon} and $n_{\min}$ in
Section~\ref{sec:methods-corruption}), so the request never exceeds the
neighbors actually available. Both resampling variants use
imbalanced-learn~\cite{lemaitre2017imbalancedlearn}.

\subsubsection{TabPFN family}

TabPFN~\cite{hollmann2023tabpfn,hollmann2025tabpfnv2} is an in-context tabular
foundation model: it is not fit with gradient descent on the study's data, but
reads the entire training set as a prompt and produces a prediction for each
test row in a single forward pass. Every TabPFN variant uses the library's
default hyperparameters, with only the random seed set per run, so no
TabPFN-specific configuration table is needed alongside
Table~\ref{tab:config}. \texttt{tabpfn\_none} takes its raw
$P(\text{minority})$ and thresholds at 0.5, and is the baseline for the rest
of the family. \texttt{tabpfn\_threshold} instead thresholds at $\pi$, the
class-1 prior used for correction, following the classical
posterior-adjustment principle of Saerens et al.~\cite{saerens2002adjusting}:
predict minority if $P(\text{minority}) > \pi$. As with every prior-dependent
method in this family, $\pi$ is either the observed, possibly corrupted
training rate $\pi_{\text{obs}}$ (the deployable variant) or the true
pre-corruption rate $\pi_{\text{true}}$ (the oracle diagnostic described in
Section~\ref{sec:methods-contribution}).

\texttt{tabpfn\_distpfn} and \texttt{tabpfn\_distpfn\_t} implement
DistPFN and DistPFN-T~\cite{lee2026distpfn}. Writing $p$ for TabPFN's raw
2-class prediction on one test row and $\bar{p}$ for the mean of $p$ across
the whole test batch,
\begin{equation}
\text{prior} = [1-\pi,\ \pi], \qquad
\text{adjusted} = \text{Normalize}\!\left(\frac{p \cdot \bar{p}}{\text{prior}}\right),
\label{eq:distpfn}
\end{equation}
using the same class-1 prior $\pi$ (observed or oracle, as defined above).
DistPFN-T adds a shared temperature
$\tau$ before the same final step:
\begin{equation}
\tau = -\sum_i \bar{p}_i \log(\text{prior}_i), \qquad
\bar{p}_T = \text{softmax}(\bar{p}/\tau),
\label{eq:distpfnt}
\end{equation}
with $\bar{p}_T$ substituted for $\bar{p}$ in Equation~\ref{eq:distpfn}. Both
formulas were verified directly against the authors' official reference
implementation before use in this study. That verification surfaced a
defect in the reference implementation itself, reported here rather than
left implicit: the authors' own published DistPFN-T reference driver
raises a shape error (\texttt{numpy}'s \texttt{AxisError}) for binary
classification, which its own exception handling silently swallows, so the
reference driver as published cannot produce a DistPFN-T score for a
two-class problem. Equation~\ref{eq:distpfnt} implements the algorithm's
evident intent, with that shape bug corrected rather than reproduced.

\texttt{tabpfn\_downsample} is the family's one data-level correction:
whichever class has more observed rows is randomly downsampled, without
replacement, to match the observed count of the smaller class, and TabPFN
reads the resulting balanced set as its prompt. Identifying the larger
class from the observed labels directly, rather than assuming a fixed
encoding stays the larger one, means this never fails regardless of how
contamination shifts the observed class balance, including at $r{=}0.50$
under Direction A, where relabeling majority rows as minority can push
the observed minority count above the observed majority count.

For prior-dependent methods (\texttt{tabpfn\_threshold},
\texttt{tabpfn\_distpfn}, \texttt{tabpfn\_distpfn\_t}), every stress cell
(one dataset, seed, ratio, contamination rate, and direction) is scored
twice: once using the observed prior $\pi_{\text{obs}}$, the only
value available in a real deployment, and once using the oracle prior
$\pi_{\text{true}}$, the clean pre-corruption rate, which a real system could
never know. The oracle variant exists only to compute the Prior Sensitivity
Gap (Section~\ref{sec:methods-contribution}) and is never treated as a
deployable result.

\subsection{Evaluation Metrics}
\label{sec:methods-metrics}

Every method, on every stress cell, is scored with ten metrics computed
against the fixed, clean test set: balanced accuracy, worst-class recall
(the lower of the two per-class recalls, which catches a collapse in one
class that balanced accuracy alone can hide), area under the
precision-recall curve (AUPRC), macro-averaged F1, minority-class
precision and recall, area under the ROC curve (AUROC), Brier
score~\cite{brier1950verification}, log loss, and expected calibration
error (ECE, ten equal-width bins)~\cite{guo2017calibration}, which
measures how far a model's predicted probabilities are from the true
observed frequencies.

All ten are computed identically for every method and condition using
scikit-learn's implementations~\cite{pedregosa2011scikitlearn}, so no
method is scored on a different basis than any other.

Not every metric receives equal treatment in Section~\ref{sec:results}.
Balanced accuracy is reported as the discrimination measure throughout.
Worst-class recall, AUPRC, and macro F1 correlate with it at 0.968, 0.781,
and 0.919 respectively (Pearson, computed across the full result set) and
largely reproduce the same ranking rather than adding independent
evidence. AUPRC's correlation is the loosest of the three, consistent with
this study's wide range of natural imbalance producing real divergence
between the two metrics, but not loose enough to promote it from a
secondary diagnostic to a headline measure on its own. Minority precision,
minority recall, and AUROC are computed as
secondary diagnostics, available for a targeted check but not used to
support a headline claim on their own. Brier score and expected
calibration error are reported separately in
Section~\ref{sec:results-reliability} because together they reveal a
divergence from the discrimination results that balanced accuracy alone
does not show (Section~\ref{sec:discussion-calibration}). Log loss is
computed identically to the other nine metrics but is not separately
reported, since its correlation with Brier score, 0.959, is comparably
high and does not add evidence beyond it.

\subsection{Contribution Metrics}
\label{sec:methods-contribution}

The ten metrics above measure how good a method is in isolation. The
following five metrics are proposed in this paper. They are defined here for
this study, not taken from prior work, and they answer the question this study
is actually about: whether a correction is worth applying once label noise is
present, relative to doing nothing at all. All five compare a corrected method against its
matching uncorrected baseline (\texttt{xgb\_none} or \texttt{tabpfn\_none})
at the \emph{same} dataset, seed, ratio, contamination rate, and direction.

\textbf{Correction Gain (CG).} $S$ is first computed per seed, then
averaged across the ten seeds within each (dataset, $r$, $c$, direction)
block, for the corrected method and its uncorrected baseline separately.
CG is the difference of those two seed-averaged scores,
\begin{equation}
\text{CG} = \bar{S}_{\text{corrected}} - \bar{S}_{\text{uncorrected}},
\label{eq:cg}
\end{equation}
sign-adjusted so that positive always means the correction helped: unchanged
for metrics where higher is better (balanced accuracy, AUPRC), and negated
for loss-type metrics where lower is better (Brier score, log loss). CG,
and every metric built from it below, therefore has one value per (method,
dataset, $r$, $c$, direction) block, not one value per stress cell; a
single seed's evaluation never counts as a unit of evidence on its own.
For the one block affected by the failure disclosed in
Section~\ref{sec:methods-implementation} (\texttt{tabpfn\_downsample} on
APS Failure, $r{=}0.10$, $c{=}0.05$, Direction A), $\bar{S}_{\text{corrected}}$
averages the nine seeds that completed rather than ten, while
$\bar{S}_{\text{uncorrected}}$ for that same block still averages all ten,
since the failure was specific to the corrected method's fit, not the
baseline's; this affects one of 416 blocks for that method and direction.

\textbf{Correction Harm Rate (CHR).} Across every (dataset, $r$, $c$,
direction) block tested for one method,
\begin{equation}
\text{CHR} = \frac{\#\{\text{blocks with CG} < 0\}}{\#\{\text{blocks tested}\}}.
\label{eq:chr}
\end{equation}
With 26 datasets, four ratios, and four nonzero contamination rates, this
is 416 blocks per method and direction (Table~\ref{tab:robustness}), not
the 4,160 individual seed-level evaluations a per-seed count would give.
CHR counts any negative seed-averaged gain as harmful, including
differences too small to distinguish from measurement noise.
Section~\ref{sec:results-rq1} also
reports a stricter variant that only counts a block as harmful once CG
falls below $-0.01$ on balanced accuracy, a sensitivity check on how much
of the headline rate reflects negligible differences rather than a
meaningful failure.

\textbf{Regret.} How large the harm is when a correction does fail, zero
otherwise:
\begin{equation}
\text{Regret} = \max(0,\ -\text{CG}).
\label{eq:regret}
\end{equation}
CHR answers how often a correction backfires. Regret answers how bad it is
when it does, so a method with low CHR but large Regret on the blocks where it
fails carries a different practical risk than one with high CHR but small
Regret.

\textbf{Break-even contamination rate ($c^*$).} For a method, dataset, and
imbalance ratio, $c^*$ is the smallest tested $c > 0$ at which the
seed-averaged CG first drops to zero or below. This marks the first noise
level a correction stops paying off, not necessarily the last, since CG is
not guaranteed to stay negative at every higher $c$ afterward
(Section~\ref{sec:threats-construct} quantifies how often it does not). If
CG never crosses zero within the tested range, $c^*$ is reported as greater
than the largest tested $c$ (0.30) rather than extrapolated beyond what was
actually measured.

\textbf{Prior Sensitivity Gap (PSG).} For the three prior-dependent TabPFN
methods,
\begin{equation}
\text{PSG} = S_{\text{oracle prior}} - S_{\text{observed prior}},
\label{eq:psg}
\end{equation}
comparing the real, deployable variant against the oracle diagnostic
described in Section~\ref{sec:methods-corrections}. A large positive PSG
means most of a correction's failure traces specifically to the prior being
wrong. A small PSG despite otherwise poor performance means the failure comes
from elsewhere.

As a robustness check on whether pooling the balanced condition
$r{=}0.50$ together with the three genuinely imbalanced ratios biases
these headline contribution metrics, Section~\ref{sec:results-r050}
recomputes CHR, mean CG, break-even share, PSG, and every statistical
test in Section~\ref{sec:methods-statistics} three times: once on the
full pooled set (as everywhere else in this paper), once restricted to
$r \in \{0.20, 0.10, 0.05\}$, and once restricted to $r{=}0.50$ alone.
All three use the identical formulas defined above, applied to the same
already-collected evaluations; no additional model fitting or scoring
was performed.

\subsection{Statistical Testing}
\label{sec:methods-statistics}

For each model family (XGBoost, TabPFN) and each corruption direction, the
methods in that family are compared using the Friedman
test~\cite{friedman1937ranks}, a non-parametric omnibus test for differences
among more than two related samples, applied to seed-averaged balanced
accuracy across every (dataset, $r$, $c$) block with a complete row for every
method in that family. For the three prior-dependent TabPFN methods, only
the observed-prior variant, the same deployable score reported everywhere
else in this study, enters the comparison. The oracle-prior rows are
excluded so that a method's column reflects its deployable score alone, not
an average of its observed and oracle variants. A family/direction
combination is tested only when at least ten complete blocks are
available; every combination in this study clears that bar by a wide
margin (416 condition-level blocks, 26 dataset-level blocks, per
family and direction), so the rule never excludes a comparison here but
is stated for reproducibility on a different data cut. When
the Friedman test is
significant at $p<0.05$, every pairwise comparison of methods within that
family and direction is tested with the Wilcoxon signed-rank
test~\cite{wilcoxon1945ranking}, and the resulting $p$-values are corrected
for multiple comparisons with Holm's step-down
procedure~\cite{holm1979sequential}.

Blocking by (dataset, $r$, $c$) treats every tested ratio and noise level
within one dataset as an independent unit of evidence, but conditions
drawn from the same dataset share its underlying feature distribution and
test set, so they are not fully independent of each other. As a more
conservative check, the same comparison is repeated with each dataset as
a single block ($n{=}26$), first averaging balanced accuracy within a
family and direction across every tested $r$ and $c$ for that dataset,
then applying the Friedman test to those 26 dataset-level values per
method, the standard practice for comparing algorithms across multiple
datasets~\cite{demsar2006statistical}. Every pairwise Wilcoxon comparison
is repeated at this same dataset-level block, Holm-corrected separately
from the (dataset, $r$, $c$) version, since a significant omnibus result
does not by itself establish that every individual pairwise difference
holds under the stricter standard. Section~\ref{sec:results-significance}
reports all versions.

\subsection{Implementation and Reproducibility}
\label{sec:methods-implementation}

Every (dataset, ratio, contamination rate, direction) combination is
repeated across all ten seeds (13, 37, 73, 5, 11, 17, 19, 23, 29, and 31),
producing one stress cell per seed, so no result depends on one particular
random split or one particular set of flipped rows. A
method failing on a given stress cell is
logged with its error rather than silently skipped, and does not stop any
other method or any later cell: one such failure occurred in the final run,
\texttt{tabpfn\_downsample} on APS Failure (seed 19, $r{=}0.10$,
$c{=}0.05$, Direction A) raised a singular-value decomposition
non-convergence error on that specific resampled training set, reproducibly
across repeated attempts, and is reported as a failure rather than
interpolated or dropped (Section~\ref{sec:results}). All other
$112{,}319$ of the $112{,}320$ planned evaluations completed.
Table~\ref{tab:config} records the exact software
versions, platform, and the frozen XGBoost configuration used throughout.
The TabPFN checkpoint used was \texttt{tabpfn-v3-classifier-v3\_default.ckpt}
(SHA-256~\texttt{d0d865d5\ldots8f3988}), and TabPFN inference used Apple's
Metal (MPS) GPU backend rather than CPU. Every library's internal thread
pool (OpenMP, BLAS) was forced to a single thread throughout the study,
because XGBoost and PyTorch, which TabPFN depends on, each bundle their own
threading runtime, and running both in the same process without this
restriction reproducibly hangs (or, in the reverse import order, segfaults).
At the dataset sizes used in this study, the resulting single-threaded
slowdown is negligible.

\begin{table}[t]
\centering
\caption{Software environment and frozen XGBoost configuration.}
\label{tab:config}
\begin{tabular}{@{}ll@{}}
\toprule
Component & Value \\
\midrule
\multicolumn{2}{@{}l}{\textit{Software versions}} \\
Python & 3.13.5 \\
Platform & macOS 26.5.1 (arm64) \\
NumPy & 2.5.2 \\
pandas & 3.0.5 \\
scikit-learn & 1.9.0 \\
SciPy & 1.18.0 \\
XGBoost & 3.4.1 \\
imbalanced-learn & 0.14.2 \\
statsmodels & 0.14.6 \\
TabPFN & 8.3.0 \\
Seeds & 13, 37, 73, 5, 11, 17, 19, 23, 29, 31 \\
\midrule
\multicolumn{2}{@{}l}{\textit{XGBoost configuration (frozen, all variants)}} \\
\texttt{n\_estimators} & 500 \\
\texttt{max\_depth} & 6 \\
\texttt{learning\_rate} & 0.05 \\
\texttt{subsample} & 0.8 \\
\texttt{colsample\_bytree} & 0.8 \\
\texttt{eval\_metric} & logloss \\
\bottomrule
\end{tabular}
\end{table}
